%%%%% MATH DEFINITIONS, shared by the body and Appendix A %%%%%
% Read by main.tex with %%%%% MATH DEFINITIONS, shared by the body and Appendix A %%%%%
% Read by main.tex with %%%%% MATH DEFINITIONS, shared by the body and Appendix A %%%%%
% Read by main.tex with %%%%% MATH DEFINITIONS, shared by the body and Appendix A %%%%%
% Read by main.tex with \input{math_commands}. Proofs use the shorthands at the end,
% so that a display is a line of mathematics and not a line of \frac's.

% --- sets and operators ---------------------------------------------------
\newcommand{\R}{\mathbb{R}}
\newcommand{\E}{\mathbb{E}}
\newcommand{\Var}{\operatorname{Var}}
\newcommand{\Cov}{\operatorname{Cov}}
\newcommand{\tr}{\operatorname{tr}}
\newcommand{\N}{\mathcal{N}}
\newcommand{\D}{\mathcal{D}}
\newcommand{\dist}{\operatorname{dist}}
\newcommand{\supp}{\operatorname{supp}}
\newcommand{\law}{\mathrm{Law}}
\newcommand{\ind}{\mathbf{1}}
\newcommand{\eps}{\varepsilon}

% --- objects of the paper -------------------------------------------------
\newcommand{\Sigmapost}{\Sigma_{\mathrm{post}}}
\newcommand{\Xhat}{\widehat{\Sigma}_X}
\newcommand{\rhohat}{\widehat{\rho}}
\newcommand{\vstar}{v_h^\star}
\newcommand{\vstarz}{v_0^\star}
\newcommand{\Covh}{\operatorname{Cov}_h}

% --- shorthands for the proofs --------------------------------------------
% label weight and kernel weight of atom #1 at bandwidth h
\newcommand{\pih}[1]{p_{#1}^{(h)}}
\newcommand{\wih}[1]{w_{#1}^{(h)}}
% label kernel K_h(y - y^i)
\newcommand{\KK}[1]{K_h(y-y^{#1})}
% per-atom collapsed field (x^i - x)/(1 - t)
\newcommand{\fat}[1]{\frac{x^{#1}-x}{1-t}}
% source density at the point that reaches x from atom #1
\newcommand{\src}[1]{\pi_0\!\Big(\frac{x-tx^{#1}}{1-t}\Big)}

% --- pointer from a statement in the body to its proof --------------------
\newcommand{\proofin}[1]{\hfill\textnormal{\footnotesize[Proof: Appendix~\ref{#1}]}}
. Proofs use the shorthands at the end,
% so that a display is a line of mathematics and not a line of \frac's.

% --- sets and operators ---------------------------------------------------
\newcommand{\R}{\mathbb{R}}
\newcommand{\E}{\mathbb{E}}
\newcommand{\Var}{\operatorname{Var}}
\newcommand{\Cov}{\operatorname{Cov}}
\newcommand{\tr}{\operatorname{tr}}
\newcommand{\N}{\mathcal{N}}
\newcommand{\D}{\mathcal{D}}
\newcommand{\dist}{\operatorname{dist}}
\newcommand{\supp}{\operatorname{supp}}
\newcommand{\law}{\mathrm{Law}}
\newcommand{\ind}{\mathbf{1}}
\newcommand{\eps}{\varepsilon}

% --- objects of the paper -------------------------------------------------
\newcommand{\Sigmapost}{\Sigma_{\mathrm{post}}}
\newcommand{\Xhat}{\widehat{\Sigma}_X}
\newcommand{\rhohat}{\widehat{\rho}}
\newcommand{\vstar}{v_h^\star}
\newcommand{\vstarz}{v_0^\star}
\newcommand{\Covh}{\operatorname{Cov}_h}

% --- shorthands for the proofs --------------------------------------------
% label weight and kernel weight of atom #1 at bandwidth h
\newcommand{\pih}[1]{p_{#1}^{(h)}}
\newcommand{\wih}[1]{w_{#1}^{(h)}}
% label kernel K_h(y - y^i)
\newcommand{\KK}[1]{K_h(y-y^{#1})}
% per-atom collapsed field (x^i - x)/(1 - t)
\newcommand{\fat}[1]{\frac{x^{#1}-x}{1-t}}
% source density at the point that reaches x from atom #1
\newcommand{\src}[1]{\pi_0\!\Big(\frac{x-tx^{#1}}{1-t}\Big)}

% --- pointer from a statement in the body to its proof --------------------
\newcommand{\proofin}[1]{\hfill\textnormal{\footnotesize[Proof: Appendix~\ref{#1}]}}
. Proofs use the shorthands at the end,
% so that a display is a line of mathematics and not a line of \frac's.

% --- sets and operators ---------------------------------------------------
\newcommand{\R}{\mathbb{R}}
\newcommand{\E}{\mathbb{E}}
\newcommand{\Var}{\operatorname{Var}}
\newcommand{\Cov}{\operatorname{Cov}}
\newcommand{\tr}{\operatorname{tr}}
\newcommand{\N}{\mathcal{N}}
\newcommand{\D}{\mathcal{D}}
\newcommand{\dist}{\operatorname{dist}}
\newcommand{\supp}{\operatorname{supp}}
\newcommand{\law}{\mathrm{Law}}
\newcommand{\ind}{\mathbf{1}}
\newcommand{\eps}{\varepsilon}

% --- objects of the paper -------------------------------------------------
\newcommand{\Sigmapost}{\Sigma_{\mathrm{post}}}
\newcommand{\Xhat}{\widehat{\Sigma}_X}
\newcommand{\rhohat}{\widehat{\rho}}
\newcommand{\vstar}{v_h^\star}
\newcommand{\vstarz}{v_0^\star}
\newcommand{\Covh}{\operatorname{Cov}_h}

% --- shorthands for the proofs --------------------------------------------
% label weight and kernel weight of atom #1 at bandwidth h
\newcommand{\pih}[1]{p_{#1}^{(h)}}
\newcommand{\wih}[1]{w_{#1}^{(h)}}
% label kernel K_h(y - y^i)
\newcommand{\KK}[1]{K_h(y-y^{#1})}
% per-atom collapsed field (x^i - x)/(1 - t)
\newcommand{\fat}[1]{\frac{x^{#1}-x}{1-t}}
% source density at the point that reaches x from atom #1
\newcommand{\src}[1]{\pi_0\!\Big(\frac{x-tx^{#1}}{1-t}\Big)}

% --- pointer from a statement in the body to its proof --------------------
\newcommand{\proofin}[1]{\hfill\textnormal{\footnotesize[Proof: Appendix~\ref{#1}]}}
. Proofs use the shorthands at the end,
% so that a display is a line of mathematics and not a line of \frac's.

% --- sets and operators ---------------------------------------------------
\newcommand{\R}{\mathbb{R}}
\newcommand{\E}{\mathbb{E}}
\newcommand{\Var}{\operatorname{Var}}
\newcommand{\Cov}{\operatorname{Cov}}
\newcommand{\tr}{\operatorname{tr}}
\newcommand{\N}{\mathcal{N}}
\newcommand{\D}{\mathcal{D}}
\newcommand{\dist}{\operatorname{dist}}
\newcommand{\supp}{\operatorname{supp}}
\newcommand{\law}{\mathrm{Law}}
\newcommand{\ind}{\mathbf{1}}
\newcommand{\eps}{\varepsilon}

% --- objects of the paper -------------------------------------------------
\newcommand{\Sigmapost}{\Sigma_{\mathrm{post}}}
\newcommand{\Xhat}{\widehat{\Sigma}_X}
\newcommand{\rhohat}{\widehat{\rho}}
\newcommand{\vstar}{v_h^\star}
\newcommand{\vstarz}{v_0^\star}
\newcommand{\Covh}{\operatorname{Cov}_h}

% --- shorthands for the proofs --------------------------------------------
% label weight and kernel weight of atom #1 at bandwidth h
\newcommand{\pih}[1]{p_{#1}^{(h)}}
\newcommand{\wih}[1]{w_{#1}^{(h)}}
% label kernel K_h(y - y^i)
\newcommand{\KK}[1]{K_h(y-y^{#1})}
% per-atom collapsed field (x^i - x)/(1 - t)
\newcommand{\fat}[1]{\frac{x^{#1}-x}{1-t}}
% source density at the point that reaches x from atom #1
\newcommand{\src}[1]{\pi_0\!\Big(\frac{x-tx^{#1}}{1-t}\Big)}

% --- pointer from a statement in the body to its proof --------------------
\newcommand{\proofin}[1]{\hfill\textnormal{\footnotesize[Proof: Appendix~\ref{#1}]}}
